%    Latex Template for ACA 2015
%%%%%%%%%%%%%%%%%%%%%%%%%%%%%%%%%%%%%%%%
%    DO NOT CHANGE THE SETTINGS BELOW
%%%%%%%%%%%%%%%%%%%%%%%%%%%%%%%%%%%%%%%%
\documentclass[11pt]{article}
\usepackage{epsfig}
\usepackage{graphicx}
\usepackage[latin1]{inputenc}
\usepackage[T1]{fontenc}
\usepackage{mathptmx}
\usepackage{url}



\begin{document}



%%%%%%%%%%%%%%%%%%%%%%%%
%%%%%%%%%%%%%%%%%%%%%%%%%%%%%%%%%%%%%%%%
%    DO NOT CHANGE THE SETTINGS ABOVE
%%%%%%%%%%%%%%%%%%%%%%%%%%%%%%%%%%%%%%%%
%
\vspace{-0.8cm}
\begin{flushleft}
\Large \textbf{\noindent
%%% PLEASE INSERT THE TITLE OF YOUR TALK HERE %%%%%%%
Dynamic Computer Illustrations and Didactic Considerations in the Learning and
Teaching of Mathematics}
%%%%%%%%%%%%%%%%%%%%%%%%%%%%%%%%%%%%%%%%%%%%%%%%%%%%%%%%
\\
%%%%%%%%%%%%%%%%%%%%%%%%%%%%%%
\vspace{0.5cm}
\normalsize
 %%%%%%             NAMES OF AUTHORS         %%%%%%
 %%%%%% NAME OF SPEAKER SHOULD BE UNDERLINED %%%%%%%%%
\normalsize{
 %%%%%% FOR EXAMPLE %%%%%
Michal Fraenkel$^1$
 %%%%%%%%%%%%%%%%%%%%%%%%
} \\
\vspace{5mm}
\textit{\footnotesize
 %%%%%% AFFILIATION OF AUTHORS %%%%%%
$^1$Center for Educational Technology, Israel
 %%%%%%%%%%%%%%%%%%%%%%%%%%%%%%%%%%%%%
}
\end{flushleft}

Years ago, as a math teacher, I used to dream of a dynamic way to show my
studen ts mathematical concepts and situations, such as rotating graphs
around a rotati on axis, graphs of functions changing according to the
change of parameters, the range of different situations meeting a certain
set of data, etc.

This is no longer a dream --- the tools are already here: We have
dynamic software that opens for us thousands of new ways to show our
students this fascinating world called ``mathematics'' --- alongside which
arise thousands of new questions.

How does the use of dynamic computer illustrations affect users' way of
thinking ? How does it affect the way teachers think? The way students
think? If using dy namic illustrations has any disadvantages, what may
they be?

In my talk, I will show various Geogebra illustrations developed for
high-school students.  I'll discuss different aspects of using them and
offer possible considerations concerning questions such as:
\begin{itemize}
\item When should we use a dynamic illustration, and when should we avoid it?
\item Should the students' age and level of the class be taken into account
when considering the use of dynamic computer illustrations?
\item What other considerations may help a teacher decide whether or not
to use a dynamic computer illustration?
\item Once a teacher decides to use a dynamic computer illustration, what
considerations should he or she take into account while actually using it
in their classroom?
\item What considerations should be taken into account while developing
dynamic computer illustrations?
\end{itemize}



%\begin{eqnarray}
%\lim_{x \to +\infty} {\rm Abstract}(x) = 1\ {\rm page}.
%\label{eq1}
%\end{eqnarray}

%%%%%%%%%%%%%%%%%%%%%%%%%%%%%%%%%%%%%%%


%%%%%%%%%%%%%%%%%%%%%%%%%%%%%%%%%%%%%%%%
%% BIBLIOGRAPHY
%%%%%%%%%%%%%%%%%%%%%%%%%%%%%%%%%%%%%%%%
%\begin{thebibliography}{00}
%\addcontentsline{toc}{chapter}{References}
%\setlength{\itemsep}{-1mm}
%\footnotesize
%\bibitem{KR} B. Keller and I. Reiten, \textit{Cluster-titled algebras are Gorenstein and stably Calabi-Yau}, Adv. Math. \textbf{211}, 1, pp. 123-151 (2007).
%\bibitem{WK} S. Washito and K. Karabashi, \textit{Use This Template}, in \textit{Proceedings of the 5th
%International Conference on Photon Interaction} (ICPI-2012), San Francisco, USA, ed. A. Chen, pp. 22-24 (2012).
%\bibitem{ABS} A.B. Smith, \textit{Describing the style for the AMMCS-CAIMS-2015 book of abstracts}, in C.D. Science (ed.), General Book of Wisdom, 2\textsuperscript{nd} ed., University of Queensland, pp. 23-42 (2011).
%\end{thebibliography}
%\normalsize
\end{document} 
